%% Automatic-driving algorithm expert resume.
%% Derived from adongwanai/LLM-Resume-Template under CC BY 4.0.

\documentclass[zh]{resume}

\fileinfo{%
  \faCopyright{} 2026, 楼振雄 \hspace{0.5em}
  \faEdit{} \today
}

\name{振雄}{楼}
\tagline{自动驾驶算法专家 \textbar{} 决策规划 · 轨迹优化}
\keywords{自动驾驶算法专家, 决策规划, 运动规划, 轨迹优化, 决策树, 凸优化, MPC}

\profile{
  \icontext{\faMapMarkerAlt}{杭州}
  \mobile{15927316272}
  \email{mrmars798@gmail.com}
  \university{华中科技大学2014~2018，2018~2021}
  \degree{机械工程 \textbullet{} 硕士}
  \info{L4港口自动驾驶经验 \textbullet{} 轨迹优化、控制算法、横向决策算法负责人}
}

\begin{document}
\makeheader

\begin{abstract}
多年港口物流自动驾驶算法研发经验，工作覆盖行为决策、实时轨迹优化、规划空间表达、车辆动力学与控制建模。长期承担规划核心模块的方案设计、代码实现、评审上线与问题闭环，近年来主导蒙特卡罗树行为搜索评估框架、实时动态避障轨迹优化器、规划语义空间重构等核心决策规划内容。方案覆盖集卡、智能导引车、物流小车等多车型和 新加坡、新西兰、国内各大港口、物流车辆等多项目。具备从问题建模、数值优化到车辆可执行性和现场交付的跨层视角。
\end{abstract}

\sectionTitle{专业技能}{\faWrench}
\begin{itemize}
  \item \textbf{决策规划}: 从零阶优化视角理解和建模决策问题，掌握搜索、采样、优化等算法，具备Monte Carlo 式随机搜索树展开/rollout框架的设计经验和多目标cost方案设计经验。
  \item \textbf{轨迹优化}: 熟悉 UTM/Frenet 参数曲线和车辆模型微分平坦表达，预研并主导QP/SQP 轨迹优化方案上线，制定多车型复杂碰撞约束、运动学约束建模方法，持续迭代轨迹优化向低耗时和高精度的目标提升。
  \item \textbf{规划空间}: 基于自动驾驶规划问题理解，全面替换原参考线的基础数据结构，大幅提升计算性能，并支持非单射投影问题的多值逆映射及物理空间回投，能够从定义层面分析参考线投影和复杂拓扑问题。
  \item \textbf{控制建模}: 具备车辆动力学、MPC、闭环模型、状态辨识及 IGV 双轴控制经验，局部控制器，观测器迭代和部署经验，对卡尔曼滤波具备理解。
  \item \textbf{工程能力}: 具有多年 C++、Linux、Git 使用经验。具备多版本上线维护、代码评审及客户现场交付经验。
\end{itemize}

\sectionTitle{工作经历}{\faBriefcase}
\begin{itemize}
  \item \textbf{杭州飞步科技股份有限公司 \quad 自动驾驶算法工程师（规划/控制）} \hfill 2021.06 --- 至今
  \begin{itemize}
    \item 个人开发: 21~22年从事车辆动力学建模，控制算法以及规控闭环分析模型；23~24年逐步负责 Path 优化器设计和维护；24~负责横向行为决策和规划空间表达等核心模块，形成决策、规划、控制与车辆建模的完整技术视角，具备问题建模的深入理解。
    \item 团队支撑: 23年开始承担带2~3人小团队的指导工作，负责或深度参与需求拆解、算法预研、框架设计、核心实现、Code Review、上线方案和长期维护。参与和指导过RRT离线寻优，动态障碍物避障，动态障碍物滤波，多车调度算法等多算法功能的开发和上线；
    \item 其他: 支撑多车型、多港区常态运营与新项目交付，参与 PSA 项目两周新加坡驻场调试并面向客户解释算法方案；22年后历次均获得团队top绩效评价。
  \end{itemize}
\end{itemize}

\sectionTitle{代表项目}{\faCode}
\begin{itemize}
  \item \textbf{决策规划重构与规划空间表达：低参考线依赖的行为搜索评估框架} \hfill 2025.06 --- 至今
  \begin{itemize}
    \item \textbf{问题背景}: 封闭场景自动驾驶方案存在强规则化路径需求，历来以地图参考线标注方案为主，强依赖地图参考线维护，选道、绕行和作业姿态调整扩展成本高。
    \item \textbf{决策方法}: 一阶段(已上线)以分治思想，在行为间给定切割点，进行变道，转弯等局部行为评估，显著减小问题求解规模，提升局部行为合理性；二阶段(推进中)将行为决策抽象为动作模版序列的零阶优化评估问题，以 Monte Carlo 式随机搜索树执行候选序列展开/rollout，设计归一化cost feature类，支持多指标、多数据源的综合归一化评估；三阶段(预研)以SL+RL方案接入Monte Carlo代理评价函数，提供更准确快速的前馈评价。
    \item \textbf{采样评估}: 根据场景和几何关系动态匹配贝塞尔行为模版，覆盖路口选道、变道绕行、泊位大范围绕行等；进行 SDF 安全距离指标、运动学可执行裕度指标和规则偏好指标，单次评估约 0.1\,ms，单行为评估维持 ms 级。
    \item \textbf{空间重构}: 提出以进度偏序、局部横向可比和多值逆映射为核心的规划坐标空间，显式表达复杂拓扑歧义并保持物理空间可回投；初步验证投影耗时下降约 90\%，构造耗时下降约 50\%，支持障碍物一对多映射。
    \item \textbf{工程结果}: 算法框架已上线至主线，稳健推进二三阶段算法迭代计划。全面支持新加坡，新西兰和国内各大港口的自动驾驶作业任务，同时支持城市场景物流车辆项目，在各场景下均具备优秀拓展性和低维护成本。
  \end{itemize}

  \item \textbf{实时 QP Path 优化器：多车型多场景统一轨迹计算层} \hfill 2023.01 --- 2025.01, 25年后负责模块维护
  \begin{itemize}
    \item \textbf{问题背景}: 业务场景中需常态化支持高精度绕行，例如港机交互，拥堵路口等，且车型覆盖带挂卡车，智能牵引车(IGV)和常规中小型车辆，需要基于不同车辆运动学规律进行适配，非线性 whole-body 优化精度高但难以满足在线时延，需要兼顾复杂车体避障精度、实时性与工程鲁棒性。
    \item \textbf{技术方案}: 主导去参考线依赖的 QP Path 优化器，以多项式曲线进行微分平坦表示，将运动学和集卡/IGV 全车碰撞约束在初始值附近线性化，并通过松弛与约束组合控制线性化残差，推导误差线性化误差上下界。
    \item \textbf{工程适配}: 在耗时优化上：引入碰撞预计算机制进行约束降采样，降低约束维度；以历史轨迹为初值做优化热启动，跨帧迭代收敛；优化矩阵前处理，稀疏矩阵优化乘法计算。在安全性上，引入轨迹复用机制，根据上一帧规划速度等信息决定车辆轨迹复用范围，保证横向稳定性和安全性。
    \item \textbf{可解释性}: 建立 debug、逐 box 碰撞分析和约束松弛诊断，将优化失败原因定位到障碍物 ID 与关键点，为决策层反馈和线上问题闭环提供依据。
    \item \textbf{最终效果}: 统一支持静态绕行、倒车、变道、靠边停车和侧向避让等工况；全工况求解保持 10\,ms 量级，典型线上统计成功率 99.88\%（350014/350439）。
  \end{itemize}

  \item \textbf{动态障碍物横向避让：横纵向协同规划能力} \hfill 2025.01 --- 2025.12
  \begin{itemize}
    \item \textbf{问题背景}: 基于横纵向耦合优化相关工作开展前期预研，推演算法在现有规划框架下的适用性与迭代求解特性，形成动态障碍物横向避让的整体方案。
    \item \textbf{技术方案}: 以交互区域相对自车路径几何特征，计算是否避让；以到交互区域的时间和障碍物加减速概率模型，计算通行让行概率指标；综合交互区域绕行倾向指标，令障碍物具备动态避让能力。
    \item \textbf{工程适配}: 全称主导方案预研和落地，开发阶段负责核心代码评审并制定测试关注项；功能稳定上线主线，为多车并行和复杂路口穿行提供可解释的横纵向联合避让能力。
  \end{itemize}

  \item \textbf{车辆动力学、控制建模与新车型适配} \hfill 2021.06 --- 2023.12, 24年后负责开发工作指导和模块维护
  \begin{itemize}
    \item 建立半挂车动力学模型并分析数值发散机理，在 20\,ms 控制周期下保持稳定计算，支持 MPC 内置模型、状态观测与仿真校验；同车型横向控制误差由上一版本的 1\,m 量级降至 0.1\,m 量级。
    \item 构建控制器闭环预测和性能分析工具，以分段线性映射解释控制输入输出边界；5\,m 前方位置预测误差由约 $\pm$0.3\,m 改善至约 $\pm$0.05\,m。
    \item 支持 IGV 前后双轴控制和轮转角零偏辨识，实现前后轴转向角自适应控制；典型全天数据横向误差不超过 0.1\,m，朝向误差不超过 0.03\,rad。
    \item 长期负责维护控制模块。参与港机对位高精度停车开发、辨识模块持续优化、新车型底盘适配等，多次驻场完成现场控制器调试。
  \end{itemize}
\end{itemize}

% \sectionTitle{教育与论文}{\faGraduationCap}
% \begin{itemize}
%   \item \textbf{华中科技大学 \quad 机械工程 \quad 硕士} \hfill 2018.09 --- 2021.06
%   \item \textbf{华中科技大学 \quad 机械工程 \quad 学士} \hfill 2014.09 --- 2018.06
%   \item \textbf{CCDC 2023}: 参与论文《Lateral Dynamic Model for Full-Size Semi-Trailer Container Trucks with Practical Evaluation》（非第一作者），参与半挂车横向动力学模型设计及宁波舟山港实车数据验证。
% \end{itemize}

\end{document}
